\section{引言}

\subsection{扫描探针显微镜的发展与分类}

扫描探针显微镜（scanning probe microscope, SPM）以一根尖锐探针在样品表面上方作栅格扫描，通过记录探针与样品之间的局域相互作用成像，其分辨率不取决于波长，而取决于探针尖端的尺寸与探针—样品间距的控制精度。1981 年，Binnig 与 Rohrer 利用金属探针与导电样品之间的隧道电流实现了表面原子分辨成像，即扫描隧道显微镜（scanning tunneling microscope, STM）\cite{binnig1982}。STM 要求样品与探针之间能够维持稳定的隧道电流，因而样品必须导电，绝缘体与多数生物样品无法直接测量。

1986 年，Binnig、Quate 与 Gerber 以探测微悬臂的形变代替隧道电流，把测量的物理量从电流换成力，提出了原子力显微镜（atomic force microscope, AFM）\cite{binnig1986}。这一改动的意义在于：探针与样品之间的相互作用力对样品的导电性没有要求，AFM 因而把 SPM 的适用范围从导体、半导体扩展到绝缘体，也使 SPM 从表面物理的实验手段变成材料、化工、生物、微电子等领域通用的形貌表征工具。

按所探测的物理量，SPM 家族可以划分为若干支：STM 探测隧道电流，反映表面的电子态密度；AFM 探测探针—样品之间的力，反映表面形貌与力学性质；横向力显微镜（lateral force microscope, LFM）探测悬臂的扭转，反映摩擦力与表面组分的差异；磁力显微镜（magnetic force microscope, MFM）与静电力显微镜（electrostatic force microscope, EFM）分别在抬起模式下探测磁性梯度与静电梯度；扫描近场光学显微镜（scanning near-field optical microscope, SNOM）通过亚波长孔径把光场限制在近场，突破衍射极限；扫描隧道谱（scanning tunneling spectroscopy, STS）在扫描的同时记录隧道电流—电压曲线，给出局域电子结构。本文的测量对象是形貌，同时记录了误差信号与横向力两路信号，后者即 LFM 意义上的摩擦力信号。

\subsection{原子力显微镜的工作模式}

AFM 的核心是一枚末端带有针尖的微悬臂。针尖与样品距离较近时，两者之间的范德华力、静电力与短程排斥力随距离急剧变化；悬臂受力的弯曲量由光杠杆法测量，一束激光聚焦在悬臂背面，反射光的位置由四象限光电探测器读出，形变量与光斑位移在一定范围内成正比。反馈回路以这个位移为输入，通过压电扫描器的 $z$ 向伸缩把位移维持在设定值，于是 $z$ 向电压的变化就记录了表面高度。

按针尖与样品的接触状态，AFM 分为三种模式。接触模式（contact mode）使针尖始终与样品接触，工作在排斥力区，反馈量是悬臂的静态偏转量；其横向分辨率高、扫描速度快，但法向力与摩擦力大，对软样品易造成损伤，也容易把表面松散颗粒推着走。轻敲模式（tapping mode）使悬臂在其共振频率附近自由振动，针尖周期性地敲击表面，反馈量是振幅的衰减；由于探针与样品只在振动周期的一小部分内接触，剪切力小，适合软样品与易损伤表面，代价是成像速度较慢。非接触模式（non-contact mode）工作在吸引力区，探针与样品始终不接触，适合研究吸附层与超软样品，但受表面水膜与污染层影响大，分辨率略低。本书实验的样品为硬度较高的聚碳酸酯光盘基片，表面结构稳定，选用接触模式。

三种模式之外，AFM 还有若干辅助成像方式：抬起模式（lift mode）在每一扫描行结束后把针尖抬高到固定高度再回扫，用于 MFM/EFM；力调制模式在扫描的同时对悬臂施加小幅高频调制，通过形变幅度的变化反映表面刚性差异；纳米加工模式则利用针尖与样品之间的强相互作用或电场在表面刻写结构\cite{manual_cspm}。

\subsection{光盘的结构、记录原理与格式演进}

光盘以聚碳酸酯（polycarbonate, PC）圆盘为基片，通过注塑把信息坑（pit）或预刻沟槽（pre-groove）复制在基片表面，再逐层覆盖功能膜。以可记录光盘 CD-R 为例，自读取面（激光入射面）向内部依次为：印刷层（标签，便于书写与识别）、保护层（紫外固化胶，防止反射层氧化与机械划伤）、反射层（金或银合金，提供高反射率）、有机染料记录层（花菁、酞菁或偶氮类染料，是记录介质）、聚碳酸酯基片（提供机械强度并承载沟槽结构）。基片上的预刻沟槽是螺旋状的连续槽，槽与槽之间的平坦区域称为岸（land）；沟槽的作用有两个：一是提供记录时激光的跟踪基准，二是界定染料涂布的边界。

CD-R 的记录过程是热化学过程。聚焦激光使染料层局部升温到数百摄氏度，染料分解、变色或与基片材料相互作用，形成光学性质与未记录区不同的标记，其反射率降低。读取时用较低功率的激光照射，光电探测器接收反射光，反射率的差异被解调为数字信号。这一点与压制式 CD-ROM 不同：后者的信息坑在注塑时就已成型在基片上，而 CD-R 的“坑”是写入后才产生的，因此 CD-R 的记录层必须是可发生不可逆变化的材料。由于染料层位于基片之上、反射层之下，记录前后标记的几何高度会发生变化，这一变化在 AFM 形貌图上表现为沿轨道的局部凸起或凹陷。

跟踪信号由沟槽本身提供。激光束在沟槽上偏离中心时，反射光的分布发生不对称变化，四象限探测器的差动信号（推挽信号）给出偏离方向与大小，伺服系统据此把光斑锁在沟槽中心。为保证跟踪信号有确定的频率基准，CD-R 的沟槽中心线并非直线，而是以固定频率左右摆动（wobble），摆动的调制携带着绝对时间信息，称为 ATIP（absolute time in pre-groove）。写入时，记录机从 ATIP 读出当前的时间码，据此确定写入位置与线速度。

各代光盘的差别主要体现在三个参数：激光波长、物镜数值孔径与道距。缩短波长、增大数值孔径可以把光斑聚得更小，从而缩小道距与最小记录符长度，提高面密度。表~\ref{tab:format} 列出主要格式的典型参数。对本实验而言，这些参数的意义在于：道距决定了 AFM 视场内可见的沟槽条数，10~$\mu$m 的视场在 CD 类光盘上应看到约 6 条、在 DVD 类上约 13 条、在蓝光上约 31 条；沟槽深度则决定了形貌图与横向力图的对比度。

\begin{table}[H]
\centering
\caption{主要光盘格式的典型参数}
\label{tab:format}
\small
\setstretch{1.0}
\begin{tabular}{@{}lccccc@{}}
\toprule
格式 & 激光波长/nm & 数值孔径 & 道距/$\mu$m & 最小坑长/$\mu$m & 单层容量 \\
\midrule
CD / CD-R & 780 & 0.45 & 1.60 & 0.83 & 650--700 MB \\
DVD / DVD-R & 650 & 0.60 & 0.74 & 0.40 & 4.7 GB \\
HD-DVD & 405 & 0.65 & 0.40 & 0.20 & 15 GB \\
Blu-ray & 405 & 0.85 & 0.32 & 0.15 & 25 GB \\
\bottomrule
\end{tabular}
\end{table}

需要说明的是，上述道距与最小坑长是格式标准规定的标称值，实际盘片存在公差；而 AFM 测得的表观周期还额外受扫描器标定的影响。本文第~\ref{sec:calib} 节将说明这两者在本次测量中如何区分。

\subsection{本实验的目的与内容}

本实验用接触模式 AFM 测量三张光盘的表面，内容包含四个方面：一是采集形貌、误差信号与横向力三路图像，观察同一表面在不同物理量上的响应；二是对形貌图作消除扫描线、面整平与对比度调整等处理，比较处理前后参数的变化；三是定量给出沟槽周期、峰谷深度与面粗糙度，并与格式标准的标称值对照，据此讨论扫描器横向标定的状态；四是识别并定量评估扫描线、条纹倾斜、探针卷积等常见伪迹，说明其对测量结果的影响量级。

